\section{Répétitions} 
%\label{sec:7:1:repetitions} 

Nous avons déjà vu quelques exemples de programmes \pascal\ dans lesquels une 
séquence de commandes du \robot\ ou d'appels de procédures se trouve répétée de 
nombreuses fois de suite. C'est le cas par exemple dans le programme {\idt FriseDeN\_1} 
dont l'exécution produit la frise ci-dessous~:

%\vspace{\baselineskip}
\vspace*{8mm}%
\noindent\includevspacePostScript[1cm]{FriseDeN}{0.95}{0.95}
\vspace*{-1cm}%

{
\parindent \parindentprog
\medskip

\carrenoir
\vspace{\neglskip}\mbox{\blc  \idt  \mcl program\idt \ FriseDeN\_1$;$
}\par 
\vspace{\neglskip}\mbox{\blc \blk\idt \mcl uses\idt \ \ldots\ $;$
}\par 
\pagebreak[3] 
\vspace{\neglskip}\mbox{\blc  \idt  \mcl begin\idt \ }\par 
\vspace{\neglskip}\mbox{\blc  \blk\idt  ag$;$}\par 
\vspace{\neglskip}\mbox{\blc  \blk\idt  av$;$\ av$;$\ pdt$;$\ vg$;$\ vd$;$\ pdt$;$\ av$;$\ av$;$\ pdt$;$\ vg$;$\ vd$;$\ pdt$;$\ }\par 
\vspace{\neglskip}\mbox{\blc  \blk\idt  av$;$\ av$;$\ pdt$;$\ vg$;$\ vd$;$\ pdt$;$\ av$;$\ av$;$\ pdt$;$\ vg$;$\ vd$;$\ pdt$;$\ }\par 
\vspace{\neglskip}\mbox{\blc  \blk\idt  av$;$\ av$;$\ pdt$;$\ vg$;$\ vd$;$\ pdt$;$\ av$;$\ av$;$\ pdt$;$\ vg$;$\ vd$;$\ pdt$;$\ }\par 
\vspace{\neglskip}\mbox{\blc  \blk\idt  av$;$\ av$;$\ pdt$;$\ vg$;$\ vd$;$\ pdt$;$\ av$;$\ av$;$\ pdt$;$\ vg$;$\ vd$;$\ pdt$;$\ }\par 
\vspace{\neglskip}\mbox{\blc  \blk\idt  av$;$\ av$;$\ pdt$;$\ vg$;$\ vd$;$\ pdt$;$\ av$;$\ av$;$\ pdt$;$\ vg$;$\ vd$;$\ pdt$;$\ }\par 
\vspace{\neglskip}\mbox{\blc  \blk\idt  av$;$\ av$;$\ pdt$;$\ vg$;$\ vd$;$\ pdt$;$\ av$;$\ av$;$\ pdt$;$\ vg$;$\ vd$;$\ pdt$;$\ }\par 
\vspace{\neglskip}\mbox{\blc  \blk\idt  av$;$\ av$;$\ pdt$;$\ vg$;$\ vd$;$\ pdt$;$\ av$;$\ av$;$\ pdt$;$\ vg$;$\ vd$;$\ pdt$;$\ }\par 
\vspace{\neglskip}\mbox{\blc  \blk\idt  av$;$\ av$;$\ pdt$;$\ vg$;$\ vd$;$\ pdt$;$\ av$;$\ av$;$\ pdt$;$\ vg$;$\ vd$;$\ pdt$;$\ }\par 
\vspace{\neglskip}\mbox{\blc  \blk\idt  st$;$}\par 
\vspace{\neglskip}\mbox{\blc  \idt  \mcl end\idt $.$}\par 

\medskip }


 
Pour éviter les répétitions successives des mêmes séquences d'ins\-truc\-tions,
nous allons étudier une nouvelle instruction du langage \pascal\ qui s'appelle la 
{\em boucle ``for''}\index{Boucle ``for''} (le mot anglais ``for'', qui signifie ``pour'', 
se prononce comme il s'écrit%
\footnote{\ Nous parlons de boucle ``for'' et non pas de boucle ``pour'' pour rester dans le
contexte du vocabulaire utilisé dans le programme.}). Voici le 
programme {\idt FriseDeN\_1} de dessin de la frise ci-dessous
réécrit en utilisant une boucle ``for''~:

%\vspace{\baselineskip}
\vspace*{8mm}%
\noindent\includevspacePostScript[1cm]{FriseDeNprime}{0.95}{0.95}
\vspace*{-1cm}%
\label{prog:FriseDeN_2}

{
\parindent \parindentprog
\medskip

\carrenoir
\vspace{\neglskip}\mbox{\blc  \idt  \mcl program\idt \ FriseDeN\_2$;$
}\par 
\vspace{\neglskip}\mbox{\blc \blk\idt \mcl uses\idt \ \ldots\ $;$
}\par 
\pagebreak[3] 
\vspace{\neglskip}\mbox{\blc  \blk\idt  \mcl var\idt \ J\ $:$\ integer$;$}\par 
\vspace{\neglskip}\mbox{\blc  \idt  \mcl begin\idt \ }\par 
\vspace{\neglskip}\mbox{\blc  \blk\idt  ag$;$}\par 
\pagebreak[3] 
\vspace{\neglskip}\mbox{\blc  \blk\idt  \mcl for\idt \ J\ $:$$=$\ 1\ \mcl to\idt \ 16\ \mcl do\idt }\par 
\vspace{\neglskip}\mbox{\blc  \blk\blk\idt  \mcl begin\idt }\par 
\vspace{\neglskip}\mbox{\blc  \blk\blk\blk\idt  av$;$\ av$;$\ pdt$;$\ vg$;$\ vd$;$\ pdt$;$\ }\par 
\vspace{\neglskip}\mbox{\blc  \blk\blk\idt  \mcl end\idt $;$}\par 
\vspace{\neglskip}\mbox{\blc  \blk\idt  st$;$}\par 
\vspace{\neglskip}\mbox{\blc  \idt  \mcl end\idt $.$}\par 

\medskip }



\noindent Juste après la ligne `{\idt {\mcl uses}\ \ldots\ $;$}' nous trouvons une {\em déclaration de 
\underline{var}\-ia\-ble}\index{Déclaration de variable}\index{Variable}
`{\idt  {\mcl var}\ J\ $:$\ integer$;$}'. Cette déclaration sert à indiquer qu'il va falloir que l'ordinateur utilise 
une case de sa mémoire comme {\em compteur}\index{Comp\-teur de bou\-cle}. Le mot anglais {\idt  integer} (prononcer 'inetédjeur')
signifie `entier'. Il indique que l'on compte avec des entiers 1, 2, \ldots. 
Nous devons donner un nom à ce compteur. C'est un identificateur que nous pouvons choisir 
comme il nous plaît (ici par exemple {\idt  J}). 
Ce compteur va servir à l'ordinateur pour compter le nombre de 
fois qu'il aura à répéter la séquence d'instructions dans la boucle ``for'', exactement comme nous pourrions le 
faire en comptant sur nos  doigts$\:$!

En français, la boucle `{\idt  {\mcl for}\ J\ $:=$\ 1\ {\mcl to}\ 16\ {\mcl do begin}\ \ldots\ {\mcl end}$;$}'
peut se traduire par  `pour J de 1 à 16 répéter début \ \ldots \ fin$;$'. Le mot anglais {\mcl to} se prononce 'tou'. Le mot
anglais {\mcl do} se prononce 'dou'. 

La séquence d'instructions à répéter s'écrit entre le {\mcl begin} (début) et le {\mcl end} (fin). Elle s'appelle le
{\em corps de la boucle}\index{Corps de boucle}. Le \index{Itération}{\em nombre d'itérations} c'est-à-dire le
nombre de fois qu'il faut répéter le corps de la boucle est indiqué 
après le mot {\mcl to}. Dans notre exemple, il faut répéter 16 fois
de suite le corps de la boucle  `{\idt  av$;$\ av$;$\ pdt$;$\ vg$;$\ vd$;$\ pdt$;$}'.

\section{Déclarations de variables entières} 
%\label{sec:7:3:decls} 

Il faut indiquer dans un programme \pascal\ que l'on utilise un ou plusieurs compteurs de boucles ``for''
par des {\em déclarations de variables}\index{Déclaration de variable} qui se 
placent dans le programme juste après les déclarations de
procédures~:
\begin{center}
\begin{tabular}{@{}|p{12.5cm}|@{}} \hline
{\em Programme \pascal~:}\\ \hline \hline
\begin{minipage}[t]{12.5cm}
{\medskip
\mbox{\blc\ \ \ {\mcl program}\hspace{0.3em}\blanc\hspace{0.3em}\idt NomDuProgramme$;$}\par 
\mbox{\blc\ \ \ \ \idt {\mcl uses}\ \ldots\ $;$}\par 
\mbox{\blc\ \ \ \ \small\rm  \ldots\ Déclarations des procédures du programme \ldots}\par 
\mbox{\blc\ \ \ \ {\mcl var}}\par 
\mbox{\blc\ \ \ \ \ \idt {\small\rm Identificateur1} $:$\ integer$;$}\par 
\mbox{\blc\ \ \ \ \ \idt {\small\rm Identificateur2} $:$\ integer$;$}\par 
\mbox{\blc\ \ \ \ \ \small\rm  \ldots}\par 
\mbox{\blc\ \ \ \mcl begin}\par 
\mbox{\blc\ \ \ \ \small\rm  \ldots\ Instructions du programme (séparées par des `$;$') \ldots}\par 
\mbox{\blc\ \ \ {\mcl end}\idt$.$}\par 
\medskip}
\end{minipage}\\ \hline 
\end{tabular}
\end{center}
\noindent Le mot anglais {\mcl var} (abréviation de `var\-ia\-ble') doit 
obligatoirement être sé\-pa\-ré de
l'iden\-ti\-fi\-ca\-teur qui le suit par un blanc \hspace{0.3em}\blanc\hspace{0.3em} ou un retour à la ligne.

\section{La boucle ``for''} 
%\label{sec:7:4:for} 

Une {\em boucle ``for''}\index{Boucle ``for''} a la forme suivante~:\\
\begin{center}
\begin{tabular}{@{}|p{12.5cm}|@{}} \hline
{\em Boucle ``for''~:}\\ \hline \hline
\begin{minipage}[t]{12.5cm}
{\medskip
\mbox{\blc\ \ \ {\mcl for}\hspace{0.3em}\blanc\hspace{0.3em}{\small\rm Identificateur}\ $:=$\ 1\hspace{0.3em}\blanc\hspace{0.3em}
{\mcl to}\hspace{0.3em}\blanc\hspace{0.3em} {\small\rm Nombre d'itérations} \hspace{0.3em}\blanc\hspace{0.3em}{\mcl do}}\par 
\mbox{\blc\ \ \ \ \mcl begin}\par 
\mbox{\blc\ \ \ \ \ \small\rm  \ldots\ Instructions du corps de la boucle (séparées par des `$;$') \ldots}\par 
\mbox{\blc\ \ \ \ {\mcl end}\idt$;$}\par 
\medskip}
\end{minipage}\\ \hline 
\end{tabular}
\end{center}
\noindent L'identificateur utilisé comme nom du compteur d'itérations  doit être déclaré comme 
désignant une variable entière. 
La  séquence d'instructions séparées par des `$;$' qui est répétée s'appelle le {\em corps de
la boucle}\index{Corps de boucle}.

\newpage
Il faut prendre garde à ne pas mettre de blanc \hspace{0.3em}\blanc\hspace{0.3em} entre 
le \hbox{deux-points~{\idt  $:$}} et le égal {\idt  $=$} du symbole {\idt  $:=$}. Par contre 
il ne faut pas oublier ce blanc \hspace{0.3em}\blanc\hspace{0.3em}
entre `{\mcl to}' et le `Nombre d'itérations' et entre le `Nombre d'itérations' et `{\mcl do}'.

Si le `Nombre d'itérations' est plus petit que 1 (0 par exemple)  ces instructions ne 
sont pas exécutées du tout. \underline{Une erreur grave est de mettre un point-}
\underline{virgule 
`{\idt $;$}' après le {\mcl do}}. Dans ce cas les instructions figurant entre le {\mcl begin} et
le {\mcl end} sont exécutées une seule fois. 

\begin{exemple}[Frise de pointes]{\rm
 La frise suivante~:\\[20mm]
\includevspacePostScript[1cm]{FriseDePointes}{0.95}{0.95}%
\\[1ex]
est produite par le programme ci-dessous~:

{
\parindent \parindentprog
\medskip

\carrenoir
\vspace{\neglskip}\mbox{\blc  \idt  \mcl program\idt \ FriseDePointes$;$
}\par 
\vspace{\neglskip}\mbox{\blc \blk\idt \mcl uses\idt \ \ldots\ $;$
}\par 
\pagebreak[3] 
\vspace{\neglskip}\mbox{\blc  \blk\idt  \mcl var\idt \ J\ $:$\ integer$;$}\par 
\vspace{\neglskip}\mbox{\blc  \idt  \mcl begin\idt \ }\par 
\vspace{\neglskip}\mbox{\blc  \blk\idt  ag$;$\ pqtd$;$\ }\par 
\pagebreak[3] 
\vspace{\neglskip}\mbox{\blc  \blk\idt  \mcl for\idt \ J\ $:$$=$\ 1\ \mcl to\idt \ 32\ \mcl do\idt }\par 
\vspace{\neglskip}\mbox{\blc  \blk\blk\idt  \mcl begin\idt }\par 
\vspace{\neglskip}\mbox{\blc  \blk\blk\blk\idt  lc$;$\ av$;$\ bc$;$\ pdt$;$\ vd$;$\ pdt$;$\ av$;$\ av$;$\ pdt$;$\ vd$;$\ }\par 
\vspace{\neglskip}\mbox{\blc  \blk\blk\idt  \mcl end\idt $;$}\par 
\vspace{\neglskip}\mbox{\blc  \blk\idt  st$;$}\par 
\vspace{\neglskip}\mbox{\blc  \idt  \mcl end\idt $.$}\par 

\medskip }


\noindent $\Box$}
\end{exemple}

\section{Utilisation de la boucle ``for'' dans une procédure} 
%\label{sec:7:5:for-proc} 

Les compteurs de boucles ``for'' utilisés dans une procédure doivent 
être déclarés dans la procédure, 
\index{Déclaration de variable}\index{Compteur de boucle}
comme suit~:
\newpage
\begin{center}
\begin{tabular}{@{}|p{12.5cm}|@{}} \hline
{\em Déclaration de procédure~:}\\ \hline \hline
\begin{minipage}[t]{12.5cm}
\medskip
\mbox{\blc\ \ \ {\mcl procedure}\hspace{0.3em}\blanc\hspace{0.3em}\idt NomDeLaProcedure$;$}\par 
\mbox{\blc\ \ \ \ {\mcl var}}\par 
\mbox{\blc\ \ \ \ \ \idt {\small\rm Identificateur1} $:$\ integer$;$}\par 
\mbox{\blc\ \ \ \ \ \idt {\small\rm Identificateur2} $:$\ integer$;$}\par 
\mbox{\blc\ \ \ \ \ \small\rm  \ldots}\par 
\mbox{\blc\ \ \ \mcl begin}\par 
\mbox{\blc\ \ \ \ \small\rm \ldots\ Instructions du corps de la procédure%
\footnote{\ Y compris les boucles ``for'' utilisant les compteurs Identificateur1,\ldots} 
(séparées par des `$;$') \ldots}\par 
\mbox{\blc\ \ \ {\mcl end}\idt$;$}\par 
\medskip
\end{minipage}\\ \hline 
\end{tabular}
\end{center}

Un même identificateur peut apparaître plusieurs fois
dans des dé\-cla\-ra\-tions de va\-ria\-bles 
`{\mcl var}\ {\small\rm Identificateur} $:$\ {\idt integer}$;$'
de procédures, à la condition expresse que ce soit dans des procédures
différentes. Ceci est illustré dans le programme
{\idt Fleche} ci-dessous où des compteurs de boucles
différents sont utilisés dans des procédures différentes ({\idt
Empennage} et {\idt Tige}) portent le même nom {\idt I}.

\begin{exemple}[Flèche]{\rm \ \\[7cm]
\includePostScript{FlecheVerticale}{0.95}{0.95}%
\vspace*{-7cm}%

{
\parindent \parindentprog
\medskip

\carrenoir
\vspace{\neglskip}\mbox{\blc  \idt  \mcl program\idt \ Fleche$;$
}\par 
\vspace{\neglskip}\mbox{\blc \blk\idt \mcl uses\idt \ \ldots\ $;$
}\par 
\medskip

\pagebreak[3] 
\vspace{\neglskip}\mbox{\blc  \blk\idt  \mcl procedure\idt \ Plume$;$}\par 
\vspace{\neglskip}\mbox{\blc  \blk\idt  \mcl begin\idt }\par 
\vspace{\neglskip}\mbox{\blc  \blk\blk\idt  av$;$\ p3htg$;$\ av$;$\ lc$;$\ p3htg$;$\ av$;$\ av$;$\ p3htg$;$\ bc$;$\ av$;$\ phtd$;$}\par 
\vspace{\neglskip}\mbox{\blc  \blk\idt  \mcl end\idt $;$}\par 
\medskip

\pagebreak[3] 
\vspace{\neglskip}\mbox{\blc  \blk\idt  \mcl procedure\idt \ Empennage$;$}\par 
\pagebreak[3] 
\vspace{\neglskip}\mbox{\blc  \blk\blk\idt  \mcl var\idt \ I\ $:$\ integer$;$}\par 
\vspace{\neglskip}\mbox{\blc  \blk\idt  \mcl begin\idt }\par 
\pagebreak[3] 
\vspace{\neglskip}\mbox{\blc  \blk\blk\idt  \mcl for\idt \ I\ $:$$=$\ 1\ \mcl to\idt \ 5\ \mcl do\idt }\par 
\vspace{\neglskip}\mbox{\blc  \blk\blk\blk\idt  \mcl begin\idt }\par 
\vspace{\neglskip}\mbox{\blc  \blk\blk\blk\blk\idt  Plume$;$}\par 
\vspace{\neglskip}\mbox{\blc  \blk\blk\blk\idt  \mcl end\idt $;$}\par 
\vspace{\neglskip}\mbox{\blc  \blk\idt  \mcl end\idt $;$}\par 
\medskip

\pagebreak[3] 
\vspace{\neglskip}\mbox{\blc  \blk\idt  \mcl procedure\idt \ Tige$;$}\par 
\pagebreak[3] 
\vspace{\neglskip}\mbox{\blc  \blk\blk\idt  \mcl var\idt \ I\ $:$\ integer$;$}\par 
\vspace{\neglskip}\mbox{\blc  \blk\idt  \mcl begin\idt }\par 
\pagebreak[3] 
\vspace{\neglskip}\mbox{\blc  \blk\blk\idt  \mcl for\idt \ I\ $:$$=$\ 1\ \mcl to\idt \ 10\ \mcl do\idt }\par 
\vspace{\neglskip}\mbox{\blc  \blk\blk\blk\idt  \mcl begin\idt }\par 
\vspace{\neglskip}\mbox{\blc  \blk\blk\blk\blk\idt  av$;$}\par 
\vspace{\neglskip}\mbox{\blc  \blk\blk\blk\idt  \mcl end\idt $;$}\par 
\vspace{\neglskip}\mbox{\blc  \blk\idt  \mcl end\idt $;$}\par 
\medskip

\pagebreak[3] 
\vspace{\neglskip}\mbox{\blc  \blk\idt  \mcl procedure\idt \ Pointe$;$}\par 
\vspace{\neglskip}\mbox{\blc  \blk\idt  \mcl begin\idt }\par 
\vspace{\neglskip}\mbox{\blc  \blk\blk\idt  pqtg$;$\ av$;$\ p3htd$;$\ av$;$\ pqtd$;$\ av$;$\ p3htd$;$\ av$;$\ }\par 
\vspace{\neglskip}\mbox{\blc  \blk\idt  \mcl end\idt $;$}\par 
\medskip

\vspace{\neglskip}\mbox{\blc  \idt  \mcl begin\idt }\par 
\vspace{\neglskip}\mbox{\blc  \blk\idt  EB$;$\ Empennage$;$\ Tige$;$\ Pointe$;$\ st$;$\ }\par 
\vspace{\neglskip}\mbox{\blc  \idt  \mcl end\idt $.$}\par 

\medskip }


\noindent $\Box$ }\end{exemple}


\section{Boucles ``for'' imbriquées (in\-dé\-pen\-dan\-tes)} 
%\label{sec:7:6:for-for-ind} 

Le corps d'une boucle ``for''~:

\mbox{\blc\ \ \ {\mcl for}\hspace{0.3em}\blanc\hspace{0.3em}{\small\rm Identificateur1}\ $:=$\ 1\hspace{0.3em}\blanc\hspace{0.3em}
{\mcl to}\hspace{0.3em}\blanc\hspace{0.3em} {\small\rm Nombre d'itérations1} \hspace{0.3em}\blanc\hspace{0.3em}{\mcl do}}\par 
\mbox{\blc\ \ \ \ \mcl begin}\par 
\mbox{\blc\ \ \ \ \ \small\rm  \ldots}\par 
\mbox{\blc\ \ \ \ {\mcl end}\idt$;$}\par 

\noindent peut à nouveau contenir des boucles ``for''~:

\mbox{\blc\ \ \ {\mcl for}\hspace{0.3em}\blanc\hspace{0.3em}{\small\rm Identificateur1}\ $:=$\ 1\hspace{0.3em}\blanc\hspace{0.3em}
{\mcl to}\hspace{0.3em}\blanc\hspace{0.3em} {\small\rm Nombre d'itérations1} \hspace{0.3em}\blanc\hspace{0.3em}{\mcl do}}\par 
\mbox{\blc\ \ \ \ \mcl begin}\par 
\mbox{\blc\ \ \ \ \ \small\rm  \ldots}\par 
\mbox{\blc\ \ \ \ \ {\mcl for}\hspace{0.3em}\blanc\hspace{0.3em}{\small\rm Identificateur2}\ $:=$\ 1\hspace{0.3em}\blanc\hspace{0.3em}
{\mcl to}\hspace{0.3em}\blanc\hspace{0.3em} {\small\rm Nombre d'itérations2} \hspace{0.3em}\blanc\hspace{0.3em}{\mcl do}}\par 
\mbox{\blc\ \ \ \ \ \ \mcl begin}\par 
\mbox{\blc\ \ \ \ \ \ \ \small\rm  \ldots}\par 
\mbox{\blc\ \ \ \ \ \ {\mcl end}\idt$;$}\par 
\mbox{\blc\ \ \ \ \ \small\rm  \ldots}\par 
\mbox{\blc\ \ \ \ {\mcl end}\idt$;$}\par 

\noindent On dit alors que les  boucles ``for'' sont {\em imbriquées}. 
La boucle utilisant le 
compteur {\small\rm Identificateur2} pour faire 
{\small\rm Nombre d'itérations2} itérations s'appelle la
 {\em boucle intérieure}. La boucle utilisant le compteur {\small\rm Identificateur1} pour faire 
{\small\rm Nombre d'itérations1} itérations s'appelle la {\em boucle extérieure}. Considérons par exemple
le dessin d'un carré~:
\newpage
\vbox to 0pt{\includevspacePostScript[20mm]{Carre1}{0.95}{0.95}}%

{
\parindent \parindentprog
\medskip

\carrenoir
\vspace{\neglskip}\mbox{\blc  \idt  \mcl program\idt \ Carre\_1$;$
}\par 
\vspace{\neglskip}\mbox{\blc \blk\idt \mcl uses\idt \ \ldots\ $;$
}\par 
\medskip

\pagebreak[3] 
\vspace{\neglskip}\mbox{\blc  \blk\idt  \mcl procedure\idt \ Cote$;$}\par 
\vspace{\neglskip}\mbox{\blc  \blk\blk\idt  \cmt  \{\ Dessiner\ un\ côté\ du\ carré\ et\ pivoter\ à\ droite\ \}\idt }\par 
\pagebreak[3] 
\vspace{\neglskip}\mbox{\blc  \blk\blk\idt  \mcl var\idt \ J\ $:$\ integer$;$}\par 
\vspace{\neglskip}\mbox{\blc  \blk\idt  \mcl begin\idt }\par 
\pagebreak[3] 
\vspace{\neglskip}\mbox{\blc  \blk\blk\idt  \mcl for\idt \ J\ $:$$=$\ 1\ \mcl to\idt \ 4\ \mcl do\idt }\par 
\vspace{\neglskip}\mbox{\blc  \blk\blk\blk\idt  \mcl begin\idt }\par 
\vspace{\neglskip}\mbox{\blc  \blk\blk\blk\blk\idt  av$;$\ }\par 
\vspace{\neglskip}\mbox{\blc  \blk\blk\blk\idt  \mcl end\idt $;$}\par 
\vspace{\neglskip}\mbox{\blc  \blk\blk\idt  pqtd$;$\ }\par 
\vspace{\neglskip}\mbox{\blc  \blk\idt  \mcl end\idt $;$}\par 
\medskip

\pagebreak[3] 
\vspace{\neglskip}\mbox{\blc  \blk\idt  \mcl procedure\idt \ Carre$;$}\par 
\vspace{\neglskip}\mbox{\blc  \blk\blk\idt  \cmt  \{\ Dessiner\ un\ carré\ \}\idt }\par 
\pagebreak[3] 
\vspace{\neglskip}\mbox{\blc  \blk\blk\idt  \mcl var\idt \ I\ $:$\ integer$;$}\par 
\vspace{\neglskip}\mbox{\blc  \blk\idt  \mcl begin\idt }\par 
\pagebreak[3] 
\vspace{\neglskip}\mbox{\blc  \blk\blk\idt  \mcl for\idt \ I\ $:$$=$\ 1\ \mcl to\idt \ 4\ \mcl do\idt }\par 
\vspace{\neglskip}\mbox{\blc  \blk\blk\blk\idt  \mcl begin\idt }\par 
\vspace{\neglskip}\mbox{\blc  \blk\blk\blk\blk\idt  Cote$;$\ }\par 
\vspace{\neglskip}\mbox{\blc  \blk\blk\blk\idt  \mcl end\idt $;$}\par 
\vspace{\neglskip}\mbox{\blc  \blk\idt  \mcl end\idt $;$}\par 
\medskip

\vspace{\neglskip}\mbox{\blc  \idt  \mcl begin\idt }\par 
\vspace{\neglskip}\mbox{\blc  \blk\idt  Carre$;$\ st$;$\ }\par 
\vspace{\neglskip}\mbox{\blc  \idt  \mcl end\idt $.$}\par 

\medskip }



En remplaçant les appels des procédures par leurs corps, on obtient
deux boucles imbriquées~:

%\includePostScript{Carre1}{0.95}{0.95}

{
\parindent \parindentprog
\medskip

\carrenoir
\vspace{\neglskip}\mbox{\blc  \idt  \mcl program\idt \ Carre\_2$;$
}\par 
\vspace{\neglskip}\mbox{\blc \blk\idt \mcl uses\idt \ \ldots\ $;$
}\par 
\vspace{\neglskip}\mbox{\blc  \blk\idt  \mcl var\idt }\par 
\vspace{\neglskip}\mbox{\blc  \blk\blk\idt  I\ $:$\ integer$;$}\par 
\vspace{\neglskip}\mbox{\blc  \blk\blk\idt  J\ $:$\ integer$;$}\par 
\vspace{\neglskip}\mbox{\blc  \idt  \mcl begin\idt }\par 
\vspace{\neglskip}\mbox{\blc  \blk\idt  \cmt  \{\ Dessiner\ un\ carré\ \}\idt }\par 
\vspace{\neglskip}\mbox{\blc  \blk\idt  \mcl for\idt \ I\ $:$$=$\ 1\ \mcl to\idt \ 4\ \mcl do\idt }\par 
\vspace{\neglskip}\mbox{\blc  \blk\blk\idt  \mcl begin\idt }\par 
\vspace{\neglskip}\mbox{\blc  \blk\blk\blk\idt  \cmt  \{\ Dessiner\ un\ côté\ du\ carré\ et\ pivoter\ à\ droite\ \}\idt }\par 
\vspace{\neglskip}\mbox{\blc  \blk\blk\blk\idt  \mcl for\idt \ J\ $:$$=$\ 1\ \mcl to\idt \ 4\ \mcl do\idt }\par 
\vspace{\neglskip}\mbox{\blc  \blk\blk\blk\blk\idt  \mcl begin\idt }\par 
\vspace{\neglskip}\mbox{\blc  \blk\blk\blk\blk\blk\idt  av$;$\ }\par 
\vspace{\neglskip}\mbox{\blc  \blk\blk\blk\blk\idt  \mcl end\idt $;$}\par 
\vspace{\neglskip}\mbox{\blc  \blk\blk\blk\idt  pqtd$;$\ }\par 
\vspace{\neglskip}\mbox{\blc  \blk\blk\idt  \mcl end\idt $;$}\par 
\vspace{\neglskip}\mbox{\blc  \blk\idt  st$;$\ }\par 
\vspace{\neglskip}\mbox{\blc  \idt  \mcl end\idt $.$}\par 

\medskip }



La boucle intérieure sert à construire un côté tandis que la boucle
ex\-té\-rieu\-re dessine les quatre côtés.

\section{Boucles ``for'' imbriquées (dépendantes)} 
%\label{sec:7:7:for-for-dep} 

Dans des boucles ``for'' imbriquées, le nombre d'itérations 
dans la boucle in\-té\-rieu\-re peut être dif\-fé\-rent à chaque fois que le corps de
la boucle extérieure est exécutée. Commençons par considérer l'exemple 
simple de segments de droite de longueurs 1, 2, 3, \ldots disposés en spirale
(le dessin est représenté ci-dessous à l'échelle $\frac{1}{2}$ avec la grille)~:\\[5cm]
\includePostScript{SpiraleCarree-1}{0.95}{0.95}%
\vspace*{-5cm}%

{
\parindent \parindentprog
\medskip

\carrenoir
\vspace{\neglskip}\mbox{\blc  \idt  \mcl program\idt \ SpiraleCarree\_1$;$
}\par 
\vspace{\neglskip}\mbox{\blc \blk\idt \mcl uses\idt \ \ldots\ $;$
}\par 
\pagebreak[3] 
\vspace{\neglskip}\mbox{\blc  \blk\idt  \mcl var\idt \ }\par 
\vspace{\neglskip}\mbox{\blc  \blk\blk\idt  I\ $:$\ integer$;$}\par 
\vspace{\neglskip}\mbox{\blc  \blk\blk\idt  J\ $:$\ integer$;$}\par 
\vspace{\neglskip}\mbox{\blc  \idt  \mcl begin\idt }\par 
\pagebreak[3] 
\vspace{\neglskip}\mbox{\blc  \blk\idt  \mcl for\idt \ I\ $:$$=$\ 1\ \mcl to\idt \ 22\ \mcl do\idt }\par 
\vspace{\neglskip}\mbox{\blc  \blk\blk\idt  \mcl begin\idt }\par 
\pagebreak[3] 
\vspace{\neglskip}\mbox{\blc  \blk\blk\blk\idt  \mcl for\idt \ J\ $:$$=$\ 1\ \mcl to\idt \ I\ \mcl do\idt }\par 
\vspace{\neglskip}\mbox{\blc  \blk\blk\blk\blk\idt  \mcl begin\idt }\par 
\vspace{\neglskip}\mbox{\blc  \blk\blk\blk\blk\blk\idt  av$;$}\par 
\vspace{\neglskip}\mbox{\blc  \blk\blk\blk\blk\idt  \mcl end\idt $;$}\par 
\vspace{\neglskip}\mbox{\blc  \blk\blk\blk\idt  pqtd$;$}\par 
\vspace{\neglskip}\mbox{\blc  \blk\blk\idt  \mcl end\idt $;$}\par 
\vspace{\neglskip}\mbox{\blc  \blk\idt  st$;$\ }\par 
\vspace{\neglskip}\mbox{\blc  \idt  \mcl end\idt $.$}\par 

\medskip }



Le nombre d'itérations dans la boucle ``for'' intérieure est égal à la valeur
du compteur {\idt I} de la boucle extérieure, c'est-à-dire successivement à 1, 2, 
3, \ldots, 22. Le programme {\idt SpiraleCarree\_1} est donc équivalent au
programme {\idt SpiraleCarree\_2} ci-dessous où les boucles ``for'' ont été
éliminées pour montrer exactement quelle est la séquence de commandes exécutée
par le \robot.

%\special{SpiraleCarree-1.ps 0.24 0.24 scale 0 -342 translate}

{
\parindent \parindentprog
\medskip

\carrenoir
\vspace{\neglskip}\mbox{\blc  \idt  \mcl program\idt \ SpiraleCarree\_2$;$
}\par 
\vspace{\neglskip}\mbox{\blc \blk\idt \mcl uses\idt \ \ldots\ $;$
}\par 
\vspace{\neglskip}\mbox{\blc  \idt  \mcl begin\idt }\par 
\vspace{\neglskip}\mbox{\blc  \blk\idt  av$;$\ pqtd$;$}\par 
\vspace{\neglskip}\mbox{\blc  \blk\idt  av$;$\ av$;$\ pqtd$;$}\par 
\vspace{\neglskip}\mbox{\blc  \blk\idt  av$;$\ av$;$\ av$;$\ pqtd$;$}\par 
\vspace{\neglskip}\mbox{\blc  \blk\idt  av$;$\ av$;$\ av$;$\ av$;$\ pqtd$;$}\par 
\vspace{\neglskip}\mbox{\blc  \blk\idt  av$;$\ av$;$\ av$;$\ av$;$\ av$;$\ pqtd$;$}\par 
\vspace{\neglskip}\mbox{\blc  \blk\idt  av$;$\ av$;$\ av$;$\ av$;$\ av$;$\ av$;$\ pqtd$;$}\par 
\vspace{\neglskip}\mbox{\blc  \blk\idt  av$;$\ av$;$\ av$;$\ av$;$\ av$;$\ av$;$\ av$;$\ pqtd$;$}\par 
\vspace{\neglskip}\mbox{\blc  \blk\idt  av$;$\ av$;$\ av$;$\ av$;$\ av$;$\ av$;$\ av$;$\ av$;$\ pqtd$;$}\par 
\vspace{\neglskip}\mbox{\blc  \blk\idt  av$;$\ av$;$\ av$;$\ av$;$\ av$;$\ av$;$\ av$;$\ av$;$\ av$;$\ pqtd$;$}\par 
\vspace{\neglskip}\mbox{\blc  \blk\idt  av$;$\ av$;$\ av$;$\ av$;$\ av$;$\ av$;$\ av$;$\ av$;$\ av$;$\ av$;$\ pqtd$;$}\par 
\vspace{\neglskip}\mbox{\blc  \blk\idt  av$;$\ av$;$\ av$;$\ av$;$\ av$;$\ av$;$\ av$;$\ av$;$\ av$;$\ av$;$\ av$;$\ pqtd$;$}\par 
\vspace{\neglskip}\mbox{\blc  \blk\idt  av$;$\ av$;$\ av$;$\ av$;$\ av$;$\ av$;$\ av$;$\ av$;$\ av$;$\ av$;$\ av$;$\ av$;$\ pqtd$;$}\par 
\vspace{\neglskip}\mbox{\blc  \blk\idt  av$;$\ av$;$\ av$;$\ av$;$\ av$;$\ av$;$\ av$;$\ av$;$\ av$;$\ av$;$\ av$;$\ av$;$\ av$;$\ pqtd$;$}\par 
\vspace{\neglskip}\mbox{\blc  \blk\idt  av$;$\ av$;$\ av$;$\ av$;$\ av$;$\ av$;$\ av$;$\ av$;$\ av$;$\ av$;$\ av$;$\ av$;$\ av$;$\ av$;$\ pqtd$;$}\par 
\vspace{\neglskip}\mbox{\blc  \blk\idt  av$;$\ av$;$\ av$;$\ av$;$\ av$;$\ av$;$\ av$;$\ av$;$\ av$;$\ av$;$\ av$;$\ av$;$\ av$;$\ av$;$\ av$;$\ pqtd$;$}\par 
\vspace{\neglskip}\mbox{\blc  \blk\idt  av$;$\ av$;$\ av$;$\ av$;$\ av$;$\ av$;$\ av$;$\ av$;$\ av$;$\ av$;$\ av$;$\ av$;$\ av$;$\ av$;$\ av$;$\ av$;$\ pqtd$;$}\par 
\vspace{\neglskip}\mbox{\blc  \blk\idt  av$;$\ av$;$\ av$;$\ av$;$\ av$;$\ av$;$\ av$;$\ av$;$\ av$;$\ av$;$\ av$;$\ av$;$\ av$;$\ av$;$\ av$;$\ av$;$\ av$;$\ pqtd$;$}\par 
\vspace{\neglskip}\mbox{\blc  \blk\idt  av$;$\ av$;$\ av$;$\ av$;$\ av$;$\ av$;$\ av$;$\ av$;$\ av$;$\ av$;$\ av$;$\ av$;$\ av$;$\ av$;$\ av$;$\ av$;$\ av$;$\ av$;$\ }\par 
\vspace{\neglskip}\mbox{\blc  \blk\idt  pqtd$;$}\par 
\vspace{\neglskip}\mbox{\blc  \blk\idt  av$;$\ av$;$\ av$;$\ av$;$\ av$;$\ av$;$\ av$;$\ av$;$\ av$;$\ av$;$\ av$;$\ av$;$\ av$;$\ av$;$\ av$;$\ av$;$\ av$;$\ av$;$\ }\par 
\vspace{\neglskip}\mbox{\blc  \blk\idt  av$;$\ pqtd$;$}\par 
\vspace{\neglskip}\mbox{\blc  \blk\idt  av$;$\ av$;$\ av$;$\ av$;$\ av$;$\ av$;$\ av$;$\ av$;$\ av$;$\ av$;$\ av$;$\ av$;$\ av$;$\ av$;$\ av$;$\ av$;$\ av$;$\ av$;$\ }\par 
\vspace{\neglskip}\mbox{\blc  \blk\idt  av$;$\ av$;$\ pqtd$;$}\par 
\vspace{\neglskip}\mbox{\blc  \blk\idt  av$;$\ av$;$\ av$;$\ av$;$\ av$;$\ av$;$\ av$;$\ av$;$\ av$;$\ av$;$\ av$;$\ av$;$\ av$;$\ av$;$\ av$;$\ av$;$\ av$;$\ av$;$\ }\par 
\vspace{\neglskip}\mbox{\blc  \blk\idt  av$;$\ av$;$\ av$;$\ pqtd$;$}\par 
\vspace{\neglskip}\mbox{\blc  \blk\idt  av$;$\ av$;$\ av$;$\ av$;$\ av$;$\ av$;$\ av$;$\ av$;$\ av$;$\ av$;$\ av$;$\ av$;$\ av$;$\ av$;$\ av$;$\ av$;$\ av$;$\ av$;$\ }\par 
\vspace{\neglskip}\mbox{\blc  \blk\idt  av$;$\ av$;$\ av$;$\ av$;$\ pqtd$;$}\par 
\vspace{\neglskip}\mbox{\blc  \blk\idt  st$;$}\par 
\vspace{\neglskip}\mbox{\blc  \idt  \mcl end\idt $.$}\par 

\medskip }



\section{Invariants de boucles ``for''} 
%\label{sec:7:8:inv}

Quand on dessine une frise comme celle-ci~:\\[9mm]
\includevspacePostScript[9mm]{FriseDeTildes}{0.95}{0.95}%
\vspace*{-9mm}%

\noindent avec une boucle ``for'', il faut que le \robot\ se trouve toujours dans 
la même direction
et dans la même position par rapport au prochain motif qui va être 
dessiné en exécutant 
le corps de la boucle. Autrement dit, à la fin du corps de la 
boucle il faut 
remettre le \robot\ dans la position de départ~:

{
\parindent \parindentprog
\medskip

\carrenoir
\vspace{\neglskip}\mbox{\blc  \idt  \mcl program\idt \ FriseDeTildes$;$
}\par 
\vspace{\neglskip}\mbox{\blc \blk\idt \mcl uses\idt \ \ldots\ $;$
}\par 
\medskip

\pagebreak[3] 
\vspace{\neglskip}\mbox{\blc  \blk\idt  \mcl procedure\idt \ Tilde$;$}\par 
\vspace{\neglskip}\mbox{\blc  \blk\blk\idt  \cmt  \{\ Dessiner\ un\ motif\ en\ forme\ de\ tilde\ \~{ }\ \}\idt }\par 
\vspace{\neglskip}\mbox{\blc  \blk\idt  \mcl begin\idt }\par 
\vspace{\neglskip}\mbox{\blc  \blk\blk\idt  phtd$;$\ av$;$\ phtd$;$\ av$;$\ pqtd$;$\ av$;$\ av$;$\ pqtg$;$\ av$;$\ phtg$;$\ av$;$}\par 
\vspace{\neglskip}\mbox{\blc  \blk\idt  \mcl end\idt $;$}\par 
\medskip

\pagebreak[3] 
\vspace{\neglskip}\mbox{\blc  \blk\idt  \mcl var\idt \ I\ $:$\ integer$;$}\par 
\medskip

\vspace{\neglskip}\mbox{\blc  \idt  \mcl begin\idt }\par 
\vspace{\neglskip}\mbox{\blc  \blk\idt  ag$;$}\par 
\pagebreak[3] 
\vspace{\neglskip}\mbox{\blc  \blk\idt  \mcl for\idt \ I\ $:$$=$\ 1\ \mcl to\idt \ 10\ \mcl do\idt }\par 
\vspace{\neglskip}\mbox{\blc  \blk\blk\idt  \mcl begin\idt }\par 
\vspace{\neglskip}\mbox{\blc  \blk\blk\blk\idt  Tilde$;$}\par 
\vspace{\neglskip}\mbox{\blc  \blk\blk\blk\idt  \cmt  \{\ Reculer\ au\ début\ du\ tilde\ suivant\ \}\idt }\par 
\vspace{\neglskip}\mbox{\blc  \blk\blk\blk\idt  lc$;$\ p3htg$;$\ av$;$\ bc$;$}\par 
\vspace{\neglskip}\mbox{\blc  \blk\blk\blk\idt  \cmt  \{\ Orienter\ le\ robot\ au\ nord\ \}\idt }\par 
\vspace{\neglskip}\mbox{\blc  \blk\blk\blk\idt  pqtd$;$\ }\par 
\vspace{\neglskip}\mbox{\blc  \blk\blk\idt  \mcl end\idt $;$}\par 
\vspace{\neglskip}\mbox{\blc  \blk\idt  st$;$}\par 
\vspace{\neglskip}\mbox{\blc  \idt  \mcl end\idt $.$}\par 

\medskip }


On dit que la position du \robot\ est un {\em invariant}\index{Invariant} de la boucle en ce sens
qu'elle doit être la même à chaque itération.

Si dans notre exemple on dessine un motif puis le suivant sans remettre entre-temps
le \robot\ dans la position de départ, on obtient le résultat suivant, dont
il faut avouer qu'il n'est pas très ressemblant à l'original~:\\[5cm]
\includePostScript{FriseDeTildesErr1}{0.95}{0.95}%
\vspace*{-5cm}

{
\parindent \parindentprog
\medskip

\carrenoir
\vspace{\neglskip}\mbox{\blc  \idt  \mcl program\idt \ FriseDeTildesErronee1$;$
}\par 
\vspace{\neglskip}\mbox{\blc \blk\idt \mcl uses\idt \ \ldots\ $;$
}\par 
\medskip

\pagebreak[3] 
\vspace{\neglskip}\mbox{\blc  \blk\idt  \mcl procedure\idt \ Tilde$;$}\par 
\vspace{\neglskip}\mbox{\blc  \blk\blk\idt  \cmt  \{\ Dessiner\ un\ motif\ en\ forme\ de\ tilde\ \~{ }\ \}\idt }\par 
\vspace{\neglskip}\mbox{\blc  \blk\idt  \mcl begin\idt }\par 
\vspace{\neglskip}\mbox{\blc  \blk\blk\idt  phtd$;$\ av$;$\ phtd$;$\ av$;$\ pqtd$;$\ }\par 
\vspace{\neglskip}\mbox{\blc  \blk\blk\idt  av$;$\ av$;$\ pqtg$;$\ av$;$\ phtg$;$\ av$;$}\par 
\vspace{\neglskip}\mbox{\blc  \blk\idt  \mcl end\idt $;$}\par 
\medskip

\pagebreak[3] 
\vspace{\neglskip}\mbox{\blc  \blk\idt  \mcl var\idt \ I\ $:$\ integer$;$}\par 
\medskip

\vspace{\neglskip}\mbox{\blc  \idt  \mcl begin\idt }\par 
\pagebreak[3] 
\vspace{\neglskip}\mbox{\blc  \blk\idt  \mcl for\idt \ I\ $:$$=$\ 1\ \mcl to\idt \ 10\ \mcl do\idt }\par 
\vspace{\neglskip}\mbox{\blc  \blk\blk\idt  \mcl begin\idt }\par 
\vspace{\neglskip}\mbox{\blc  \blk\blk\blk\idt  Tilde$;$}\par 
\vspace{\neglskip}\mbox{\blc  \blk\blk\idt  \mcl end\idt $;$}\par 
\vspace{\neglskip}\mbox{\blc  \blk\idt  st$;$}\par 
\vspace{\neglskip}\mbox{\blc  \idt  \mcl end\idt $.$}\par 

\medskip }



Si l'on pense à faire reculer le \robot\ mais que l'on oublie de le remettre
dans la bonne direction, le résultat est encore plus étrange~:\\[25mm]
\includePostScript{FriseDeTildesErr2}{0.95}{0.95}%
\vspace*{-25mm}

{
\parindent \parindentprog
\medskip

\carrenoir
\vspace{\neglskip}\mbox{\blc  \idt  \mcl program\idt \ FriseDeTildesErronee2$;$
}\par 
\vspace{\neglskip}\mbox{\blc \blk\idt \mcl uses\idt \ \ldots\ $;$
}\par 
\medskip

\pagebreak[3] 
\vspace{\neglskip}\mbox{\blc  \blk\idt  \mcl procedure\idt \ Tilde$;$}\par 
\vspace{\neglskip}\mbox{\blc  \blk\blk\idt  \cmt  \{\ Dessiner\ un\ motif\ en\ forme\ de\ tilde\ \~{ }\ \}\idt }\par 
\vspace{\neglskip}\mbox{\blc  \blk\idt  \mcl begin\idt }\par 
\vspace{\neglskip}\mbox{\blc  \blk\blk\idt  phtd$;$\ av$;$\ phtd$;$\ av$;$\ pqtd$;$\ av$;$\ av$;$\ pqtg$;$\ av$;$\ phtg$;$\ av$;$}\par 
\vspace{\neglskip}\mbox{\blc  \blk\idt  \mcl end\idt $;$}\par 
\medskip

\pagebreak[3] 
\vspace{\neglskip}\mbox{\blc  \blk\idt  \mcl var\idt \ I\ $:$\ integer$;$}\par 
\medskip

\vspace{\neglskip}\mbox{\blc  \idt  \mcl begin\idt }\par 
\pagebreak[3] 
\vspace{\neglskip}\mbox{\blc  \blk\idt  \mcl for\idt \ I\ $:$$=$\ 1\ \mcl to\idt \ 13\ \mcl do\idt }\par 
\vspace{\neglskip}\mbox{\blc  \blk\blk\idt  \mcl begin\idt }\par 
\vspace{\neglskip}\mbox{\blc  \blk\blk\blk\idt  Tilde$;$\ \cmt  \{\ Reculer\ au\ début\ du\ tilde\ suivant\ \}\idt \ lc$;$\ p3htg$;$\ av$;$\ bc$;$}\par 
\vspace{\neglskip}\mbox{\blc  \blk\blk\idt  \mcl end\idt $;$}\par 
\vspace{\neglskip}\mbox{\blc  \blk\idt  st$;$}\par 
\vspace{\neglskip}\mbox{\blc  \idt  \mcl end\idt $.$}\par 

\medskip }



Une bonne façon de trouver les erreurs dans une boucle ``for'' est donc
de chercher l'invariant, c'est-à-dire la position et l'orientation dans
lesquelles doit se trouver le \robot\ avant de dessiner le motif en 
exécutant le corps de la boucle. Il faut alors vérifier que le \robot\ est
convenablement placé quand on arrive pour la première fois dans la boucle
et qu'il l'est à nouveau après avoir exécuté le corps de la boucle.

La même notion d'invariant s'applique pour les procédures. Avant chaque
appel d'une procédure, il faut que le robot soit dans un ou plusieurs
états possibles qui doivent être parfaitement définis par la position 
levée ou baissée du crayon du \robot\ et par la position et l'orientation
du \robot\ par rapport à la figure dessinée par la procédure.

\noindent\rule{137mm}{0.1mm} %---------------------------------------------
\begin{exercices}{\ \\

\vspace{-\baselineskip}
\tiret Écrire des programmes \pascal\ utilisant des boucles ``for'' pour dessiner 
quelques-unes des frises de la figure 
\ref{fig:Frise-1}  page \pageref{fig:Frise-1}.

\tiret Écrire des programmes \pascal\ utilisant des boucles ``for'' dans des pro\-cé\-du\-res
pour dessiner quelques-unes des frises de la figure 
\ref{fig:Frise-2}  page \pageref{fig:Frise-2}.

\tiret Écrire des programmes \pascal\ pour dessiner quelques-unes des rosaces,
frises et pavage de la figure 
\ref{fig:Ro-Fr-Pa}  page \pageref{fig:Ro-Fr-Pa}, en utilisant des boucles
``for'' imbriquées indépendantes.

\tiret Écrire des programmes \pascal\ pour dessiner quelques-unes des spirales,
quadrillages et autres tuyaux d'orgues de la figure 
\ref{fig:SpirQuadTuy}  ci-dessous, en utilisant des boucles
``for'' imbriquées dépendantes.

\tiret Écrire un programme \pascal\ pour dessiner le pavage de 
la figure \ref{fig:PavCLC} ci-dessous en indiquant avec précision les
invariants des procédures et boucles ``for''.
\noindent $\Box$}\end{exercices}
\begin{figure}[htbp]
 \includevspacePostScript[4.5cm]{SpirQuadTuy}{0.95}{0.95}
 \caption{Quadrillages, carrés, spirale et tuyaux d'orgues}
 \label{fig:SpirQuadTuy}
\end{figure}
\begin{figure}[htbp]
 \includevspacePostScript[7.1cm]{PavageCarresTrianglesCercles}{0.95}{0.95}
 \caption{Pavage de carrés, triangles et cercles}\label{fig:PavCLC}
\end{figure}
\begin{figure}[htbp]
 \includevspacePostScript[17.8cm]{Frise-1}{0.95}{0.95}
 \caption{Frises (1)}\label{fig:Frise-1}
\end{figure}
\begin{figure}[htbp]
 \includevspacePostScript[17.8cm]{Frise-2}{0.95}{0.95}
 \caption{Frises (2)}\label{fig:Frise-2}
\end{figure}
\begin{figure}[htbp]
 \includevspacePostScript[17.8cm]{Ro-Fr-Pa}{0.95}{0.95}
 \caption{Rosaces, frises et pavage}\label{fig:Ro-Fr-Pa}
\end{figure}

\clearpage

%%%%%%%%%%%%%%%%%%%%%%%%%%%%%%%%%%%%%%%%%%%%%%%%%%%%%%%%%%%%%%%%%%%%%%%%%%%%%%%%%%%%%%%%%%%%%%
\begin{corriges}{
\renewcommand{\carrenoir}{\hspace{-4mm}\rule{1.5mm}{1.5mm}\hspace{2.5mm}}
Nous donnons quelques programmes, les autres se trouvant sur la disquette
 d'accompagnement du livre~:

{
\parindent \parindentprog
\medskip

\carrenoir
\vspace{\neglskip}\mbox{\blc  \idt  \mcl program\idt \ FriseCrantee$;$
}\par 
\vspace{\neglskip}\mbox{\blc \blk\idt \mcl uses\idt \ \ldots\ $;$
}\par 
\medskip

\pagebreak[3] 
\vspace{\neglskip}\mbox{\blc  \blk\idt  \mcl procedure\idt \ Cran$;$}\par 
\vspace{\neglskip}\mbox{\blc  \blk\blk\idt  \cmt  \{\ Dessiner\ le\ motif\ {'}n\_{'}\ \}\idt }\par 
\vspace{\neglskip}\mbox{\blc  \blk\idt  \mcl begin\idt }\par 
\vspace{\neglskip}\mbox{\blc  \blk\blk\blk\idt  av$;$\ pqtd$;$\ av$;$\ pqtd$;$\ av$;$\ pqtg$;$\ av$;$\ pqtg$;$}\par 
\vspace{\neglskip}\mbox{\blc  \blk\idt  \mcl end\idt $;$}\par 
\medskip

\pagebreak[3] 
\vspace{\neglskip}\mbox{\blc  \blk\idt  \mcl var\idt \ I\ $:$\ integer$;$}\par 
\medskip

\vspace{\neglskip}\mbox{\blc  \idt  \mcl begin\idt }\par 
\vspace{\neglskip}\mbox{\blc  \blk\idt  ag$;$\ }\par 
\pagebreak[3] 
\vspace{\neglskip}\mbox{\blc  \blk\idt  \mcl for\idt \ I$:$$=$\ 1\ \mcl to\idt \ 21\ \mcl do\idt }\par 
\vspace{\neglskip}\mbox{\blc  \blk\blk\idt  \mcl begin\idt \ }\par 
\vspace{\neglskip}\mbox{\blc  \blk\blk\blk\idt  Cran$;$}\par 
\vspace{\neglskip}\mbox{\blc  \blk\blk\idt  \mcl end\idt $;$}\par 
\vspace{\neglskip}\mbox{\blc  \blk\idt  st$;$}\par 
\vspace{\neglskip}\mbox{\blc  \idt  \mcl end\idt $.$}\par 

\medskip }



{
\parindent \parindentprog
\medskip

\carrenoir
\vspace{\neglskip}\mbox{\blc  \idt  \mcl program\idt \ QuadrillageTriangulaire$;$
}\par 
\vspace{\neglskip}\mbox{\blc \blk\idt \mcl uses\idt \ \ldots\ $;$
}\par 
\medskip

\pagebreak[3] 
\vspace{\neglskip}\mbox{\blc  \idt  \mcl var\idt }\par 
\vspace{\neglskip}\mbox{\blc  \blk\idt  I\ $:$\ integer$;$\ }\par 
\vspace{\neglskip}\mbox{\blc  \blk\idt  J\ $:$\ integer$;$}\par 
\vspace{\neglskip}\mbox{\blc  \blk\idt  K\ $:$\ integer$;$}\par 
\vspace{\neglskip}\mbox{\blc  \idt  \mcl begin\idt }\par 
\vspace{\neglskip}\mbox{\blc  \blk\idt  \cmt  \{\ Dessiner\ 10\ lignes\ de\ carrés\ juxtaposés\ horizontalement\ \}\idt }\par 
\vspace{\neglskip}\mbox{\blc  \blk\idt  pqtd$;$\ }\par 
\pagebreak[3] 
\vspace{\neglskip}\mbox{\blc  \blk\idt  \mcl for\idt \ I\ $:$$=$\ 1\ \mcl to\idt \ 10\ \mcl do\idt }\par 
\vspace{\neglskip}\mbox{\blc  \blk\blk\idt  \mcl begin\idt }\par 
\vspace{\neglskip}\mbox{\blc  \blk\blk\blk\idt  \cmt  \{\ Dessiner\ une\ ligne\ de\ \ I\ \ carrés\ juxtaposés\ horizontalement,\ \}\idt }\par 
\vspace{\neglskip}\mbox{\blc  \blk\blk\blk\idt  \cmt  \{\ en\ partant\ en\ haut\ à\ gauche\ et\ en\ arrivant\ en\ haut\ à\ droite,\ \}\idt }\par 
\vspace{\neglskip}\mbox{\blc  \blk\blk\blk\idt  \cmt  \{\ le\ robot\ étant\ dans\ les\ deux\ cas\ orienté\ vers\ l{'}est.\ \ \ \ \ \ \ \ \ \}\idt }\par 
\pagebreak[3] 
\vspace{\neglskip}\mbox{\blc  \blk\blk\blk\idt  \mcl for\idt \ J\ $:$$=$\ 1\ \mcl to\idt \ I\ \mcl do\idt }\par 
\vspace{\neglskip}\mbox{\blc  \blk\blk\blk\blk\idt  \mcl begin\idt }\par 
\vspace{\neglskip}\mbox{\blc  \blk\blk\blk\blk\blk\idt  \cmt  \{\ Dessiner\ un\ carré,\ en\ partant\ du\ coin\ supérieur\ gauche\ et\ en\ \}\idt }\par 
\vspace{\neglskip}\mbox{\blc  \blk\blk\blk\blk\blk\idt  \cmt  \{\ arrivant\ \ au\ \ coin\ \ supérieur\ droit,\ le\ robot\ étant\ dans\ les\ \}\idt }\par 
\vspace{\neglskip}\mbox{\blc  \blk\blk\blk\blk\blk\idt  \cmt  \{\ deux\ cas\ orienté\ vers\ l{'}est.\ \ \ \ \ \ \ \ \ \ \ \ \ \ \ \ \ \ \ \ \ \ \ \ \ \ \ \ \ \ \ \ \ \}\idt }\par 
\pagebreak[3] 
\vspace{\neglskip}\mbox{\blc  \blk\blk\blk\blk\blk\idt  \mcl for\idt \ K\ $:$$=$\ 1\ \mcl to\idt \ 4\ \mcl do\idt }\par 
\vspace{\neglskip}\mbox{\blc  \blk\blk\blk\blk\blk\blk\idt  \mcl begin\idt }\par 
\vspace{\neglskip}\mbox{\blc  \blk\blk\blk\blk\blk\blk\blk\idt  av$;$\ pqtd$;$\ }\par 
\vspace{\neglskip}\mbox{\blc  \blk\blk\blk\blk\blk\blk\idt  \mcl end\idt $;$}\par 
\vspace{\neglskip}\mbox{\blc  \blk\blk\blk\blk\blk\idt  av$;$}\par 
\vspace{\neglskip}\mbox{\blc  \blk\blk\blk\blk\idt  \mcl end\idt $;$}\par 
\vspace{\neglskip}\mbox{\blc  \blk\blk\blk\idt  \cmt  \{\ Revenir,\ sans\ dessiner,\ au\ début\ de\ la\ ligne\ suivante,\ le\ robot\ \}\idt }\par 
\vspace{\neglskip}\mbox{\blc  \blk\blk\blk\idt  \cmt  \{\ étant\ orienté\ vers\ l{'}est\ au\ départ\ et\ à\ l{'}arrivée.\ \ \ \ \ \ \ \ \ \ \ \ \ \ \}\idt }\par 
\vspace{\neglskip}\mbox{\blc  \blk\blk\blk\idt  lc$;$\ pdt$;$\ }\par 
\pagebreak[3] 
\vspace{\neglskip}\mbox{\blc  \blk\blk\blk\idt  \mcl for\idt \ J\ $:$$=$\ 1\ \mcl to\idt \ I\ \mcl do\idt }\par 
\vspace{\neglskip}\mbox{\blc  \blk\blk\blk\blk\idt  \mcl begin\idt }\par 
\vspace{\neglskip}\mbox{\blc  \blk\blk\blk\blk\blk\idt  av$;$}\par 
\vspace{\neglskip}\mbox{\blc  \blk\blk\blk\blk\idt  \mcl end\idt $;$}\par 
\vspace{\neglskip}\mbox{\blc  \blk\blk\blk\idt  pqtg$;$\ av$;$\ pqtg$;$\ bc$;$\ }\par 
\vspace{\neglskip}\mbox{\blc  \blk\blk\idt  \mcl end\idt $;$}\par 
\vspace{\neglskip}\mbox{\blc  \blk\idt  st$;$}\par 
\vspace{\neglskip}\mbox{\blc  \idt  \mcl end\idt $.$}\par 

\medskip }



{
\parindent \parindentprog
\medskip

\carrenoir
\vspace{\neglskip}\mbox{\blc  \idt  \mcl program\idt \ PavageCarresTrianglesCercles$;$
}\par 
\vspace{\neglskip}\mbox{\blc \blk\idt \mcl uses\idt \ \ldots\ $;$
}\par 
\medskip

\pagebreak[3] 
\vspace{\neglskip}\mbox{\blc  \blk\idt  \mcl procedure\idt \ Carre$;$}\par 
\vspace{\neglskip}\mbox{\blc  \blk\blk\idt  \cmt  \{\ Dessiner\ un\ carr\'{e}\ de\ c\^{o}t\'{e}\ de\ taille\ 2,\ le\ robot\ partant\ du\ coin\ \}\idt }\par 
\vspace{\neglskip}\mbox{\blc  \blk\blk\idt  \cmt  \{\ inf\'{e}rieur\ \ gauche\ tourn\'{e}\ vers\ l{'}est,\ et\ arrivant\ crayon\ baiss\'{e}\ \ \}\idt }\par 
\vspace{\neglskip}\mbox{\blc  \blk\blk\idt  \cmt  \{\ au\ coin\ inf\'{e}rieur\ droit\ \'{e}galement\ orient\'{e}\ \`{a}\ l{'}est.\ \ \ \ \ \ \ \ \ \ \ \ \ \ \}\idt }\par 
\vspace{\neglskip}\mbox{\blc  \blk\idt  \mcl begin\idt }\par 
\vspace{\neglskip}\mbox{\blc  \blk\blk\idt  \cmt  \{\ Dessiner\ le\ carr\'{e}.\ \ \ \ \ \ \ \ \ \ \ \ \ \ \ \ \ \ \ \ \ \ \ \ \ \ \ \ \ \ \ \ \ \ \ \ \ \ \ \ \ \ \ \ \ \ \}\idt }\par 
\vspace{\neglskip}\mbox{\blc  \blk\blk\idt  av$;$\ av$;$\ pqtg$;$\ av$;$\ av$;$\ pqtg$;$\ av$;$\ av$;$\ pqtg$;$\ av$;$\ av$;$\ pqtg$;$}\par 
\vspace{\neglskip}\mbox{\blc  \blk\blk\idt  \cmt  \{\ Avancer\ au\ coin\ inf\'{e}rieur\ droit.\ \ \ \ \ \ \ \ \ \ \ \ \ \ \ \ \ \ \ \ \ \ \ \ \ \ \ \ \ \ \ \ \}\idt }\par 
\vspace{\neglskip}\mbox{\blc  \blk\blk\idt  av$;$\ av$;$\ }\par 
\vspace{\neglskip}\mbox{\blc  \blk\idt  \mcl end\idt $;$}\par 
\medskip

\pagebreak[3] 
\vspace{\neglskip}\mbox{\blc  \blk\idt  \mcl procedure\idt \ Triangles$;$}\par 
\vspace{\neglskip}\mbox{\blc  \blk\blk\idt  \cmt  \{\ Dessiner\ \ deux\ triangles\ oppos\'{e}s\ par\ un\ sommet\ inscrits\ dans\ un\ \}\idt }\par 
\vspace{\neglskip}\mbox{\blc  \blk\blk\idt  \cmt  \{\ carr\'{e}\ invisible\ du\ quadrillage\ de\ c\^{o}t\'{e}\ 2,\ en\ \ partant\ \ du\ \ coin\ \}\idt }\par 
\vspace{\neglskip}\mbox{\blc  \blk\blk\idt  \cmt  \{\ inf\'{e}rieur\ \ gauche\ du\ carr\'{e}\ tourn\'{e}\ vers\ l{'}est,\ et\ en\ arrivant\ au\ \}\idt }\par 
\vspace{\neglskip}\mbox{\blc  \blk\blk\idt  \cmt  \{\ coin\ inf\'{e}rieur\ droit,\ crayon\ baiss\'{e}\ et\ tourn\'{e}\ vers\ l{'}est.\ \ \ \ \ \ \ \}\idt }\par 
\vspace{\neglskip}\mbox{\blc  \blk\idt  \mcl begin\idt }\par 
\vspace{\neglskip}\mbox{\blc  \blk\blk\idt  lc$;$\ av$;$\ bc$;$}\par 
\vspace{\neglskip}\mbox{\blc  \blk\blk\idt  pqtg$;$\ av$;$\ av$;$\ p3htg$;$\ av$;$\ p3htg$;$\ av$;$\ av$;$\ p3htd$;$\ av$;$\ p3htg$;$}\par 
\vspace{\neglskip}\mbox{\blc  \blk\blk\idt  lc$;$\ av$;$\ bc$;$\ }\par 
\vspace{\neglskip}\mbox{\blc  \blk\idt  \mcl end\idt $;$}\par 
\medskip

\pagebreak[3] 
\vspace{\neglskip}\mbox{\blc  \blk\idt  \mcl procedure\idt \ Cercle$;$}\par 
\vspace{\neglskip}\mbox{\blc  \blk\blk\idt  \cmt  \{\ Dessiner\ \ un\ \ cercle\ de\ rayon\ 1\ inscrit\ dans\ un\ carr\'{e}\ invisible\ \}\idt }\par 
\vspace{\neglskip}\mbox{\blc  \blk\blk\idt  \cmt  \{\ du\ quadrillage\ de\ c\^{o}t\'{e}\ 2,\ en\ partant\ du\ coin\ \ inf\'{e}rieur\ \ gauche\ \}\idt }\par 
\vspace{\neglskip}\mbox{\blc  \blk\blk\idt  \cmt  \{\ du\ \ carr\'{e}\ \ tourn\'{e}\ \ vers\ l{'}est,\ et\ en\ arrivant\ au\ coin\ inf\'{e}rieur\ \}\idt }\par 
\vspace{\neglskip}\mbox{\blc  \blk\blk\idt  \cmt  \{\ droit\ du\ carr\'{e},\ crayon\ baiss\'{e}\ et\ tourn\'{e}\ vers\ l{'}est.\ \ \ \ \ \ \ \ \ \ \ \ \ \}\idt }\par 
\vspace{\neglskip}\mbox{\blc  \blk\idt  \mcl begin\idt \ }\par 
\vspace{\neglskip}\mbox{\blc  \blk\blk\idt  lc$;$\ av$;$\ bc$;$\ }\par 
\vspace{\neglskip}\mbox{\blc  \blk\blk\idt  vg$;$\ vg$;$\ vg$;$\ vg$;$\ }\par 
\vspace{\neglskip}\mbox{\blc  \blk\blk\idt  lc$;$\ av$;$\ bc$;$\ }\par 
\vspace{\neglskip}\mbox{\blc  \blk\idt  \mcl end\idt $;$}\par 
\medskip

\pagebreak[3] 
\vspace{\neglskip}\mbox{\blc  \blk\idt  \mcl procedure\idt \ LigneCarre$;$}\par 
\vspace{\neglskip}\mbox{\blc  \blk\idt  \cmt  \{\ Dessiner\ une\ ligne\ constitu\'{e}e\ d{'}un\ motif\ r\'{e}p\'{e}t\'{e}\ 7\ fois,\ le\ motif\ \}\idt }\par 
\vspace{\neglskip}\mbox{\blc  \blk\idt  \cmt  \{\ \'{e}tant\ \ form\'{e}\ d{'}un\ carr\'{e},\ de\ triangles\ et\ d{'}un\ cercle\ juxtapos\'{e}s,\ \}\idt }\par 
\vspace{\neglskip}\mbox{\blc  \blk\idt  \cmt  \{\ le\ robot\ partant\ du\ coin\ inf\'{e}rieur\ gauche\ de\ \ la\ \ ligne\ \ orient\'{e}\ \}\idt }\par 
\vspace{\neglskip}\mbox{\blc  \blk\idt  \cmt  \{\ vers\ l{'}est,\ et\ arrivant\ au\ coin\ inf\'{e}rieur\ droit\ \'{e}galement\ orien-\ \}\idt }\par 
\vspace{\neglskip}\mbox{\blc  \blk\idt  \cmt  \{\ t\'{e}\ \`{a}\ l{'}est.\ \ \ \ \ \ \ \ \ \ \ \ \ \ \ \ \ \ \ \ \ \ \ \ \ \ \ \ \ \ \ \ \ \ \ \ \ \ \ \ \ \ \ \ \ \ \ \ \ \ \ \ \ \ \}\idt }\par 
\pagebreak[3] 
\vspace{\neglskip}\mbox{\blc  \blk\blk\idt  \mcl var\idt \ I\ $:$\ integer$;$}\par 
\vspace{\neglskip}\mbox{\blc  \blk\idt  \mcl begin\idt }\par 
\pagebreak[3] 
\vspace{\neglskip}\mbox{\blc  \blk\blk\idt  \mcl for\idt \ I\ $:$$=$\ 1\ \mcl to\idt \ 7\ \mcl do\idt }\par 
\vspace{\neglskip}\mbox{\blc  \blk\blk\blk\idt  \mcl begin\idt }\par 
\vspace{\neglskip}\mbox{\blc  \blk\blk\blk\blk\idt  \cmt  \{\ Le\ robot\ est\ crayon\ baiss\'{e},\ \`{a}\ l{'}est,\ en\ bas\ \`{a}\ gauche\ du\ motif\ \}\idt }\par 
\vspace{\neglskip}\mbox{\blc  \blk\blk\blk\blk\idt  Carre$;$\ Triangles$;$\ Cercle$;$\ }\par 
\vspace{\neglskip}\mbox{\blc  \blk\blk\blk\idt  \mcl end\idt $;$}\par 
\vspace{\neglskip}\mbox{\blc  \blk\idt  \mcl end\idt $;$}\par 
\medskip

\pagebreak[3] 
\vspace{\neglskip}\mbox{\blc  \blk\idt  \mcl procedure\idt \ LigneTriangles$;$}\par 
\vspace{\neglskip}\mbox{\blc  \blk\idt  \cmt  \{\ Dessiner\ une\ ligne\ constitu\'{e}e\ d{'}un\ motif\ r\'{e}p\'{e}t\'{e}\ 7\ fois,\ le\ motif\ \}\idt }\par 
\vspace{\neglskip}\mbox{\blc  \blk\idt  \cmt  \{\ \'{e}tant\ \ form\'{e}\ de\ triangles,\ d{'}un\ cercle\ et\ d{'}un\ carr\'{e}\ juxtapos\'{e}s,\ \}\idt }\par 
\vspace{\neglskip}\mbox{\blc  \blk\idt  \cmt  \{\ le\ robot\ partant\ du\ coin\ inf\'{e}rieur\ gauche\ de\ \ la\ \ ligne\ \ orient\'{e}\ \}\idt }\par 
\vspace{\neglskip}\mbox{\blc  \blk\idt  \cmt  \{\ vers\ l{'}est,\ et\ arrivant\ au\ coin\ inf\'{e}rieur\ droit\ \'{e}galement\ orien-\ \}\idt }\par 
\vspace{\neglskip}\mbox{\blc  \blk\idt  \cmt  \{\ t\'{e}\ \`{a}\ l{'}est.\ \ \ \ \ \ \ \ \ \ \ \ \ \ \ \ \ \ \ \ \ \ \ \ \ \ \ \ \ \ \ \ \ \ \ \ \ \ \ \ \ \ \ \ \ \ \ \ \ \ \ \ \ \ \}\idt }\par 
\pagebreak[3] 
\vspace{\neglskip}\mbox{\blc  \blk\blk\idt  \mcl var\idt \ I\ $:$\ integer$;$}\par 
\vspace{\neglskip}\mbox{\blc  \blk\idt  \mcl begin\idt }\par 
\pagebreak[3] 
\vspace{\neglskip}\mbox{\blc  \blk\blk\idt  \mcl for\idt \ I\ $:$$=$\ 1\ \mcl to\idt \ 7\ \mcl do\idt }\par 
\vspace{\neglskip}\mbox{\blc  \blk\blk\blk\idt  \mcl begin\idt }\par 
\vspace{\neglskip}\mbox{\blc  \blk\blk\blk\blk\idt  \cmt  \{\ Le\ robot\ est\ crayon\ baiss\'{e},\ \`{a}\ l{'}est,\ en\ bas\ \`{a}\ gauche\ du\ motif\ \}\idt }\par 
\vspace{\neglskip}\mbox{\blc  \blk\blk\blk\blk\idt  Triangles$;$\ Cercle$;$\ Carre$;$\ }\par 
\vspace{\neglskip}\mbox{\blc  \blk\blk\blk\idt  \mcl end\idt $;$}\par 
\vspace{\neglskip}\mbox{\blc  \blk\idt  \mcl end\idt $;$}\par 
\medskip

\pagebreak[3] 
\vspace{\neglskip}\mbox{\blc  \blk\idt  \mcl procedure\idt \ LigneCercle$;$}\par 
\vspace{\neglskip}\mbox{\blc  \blk\idt  \cmt  \{\ Dessiner\ une\ ligne\ constitu\'{e}e\ d{'}un\ motif\ r\'{e}p\'{e}t\'{e}\ 7\ fois,\ le\ motif\ \}\idt }\par 
\vspace{\neglskip}\mbox{\blc  \blk\idt  \cmt  \{\ \'{e}tant\ \ form\'{e}\ d{'}un\ cercle,\ de\ triangles\ et\ d{'}un\ carr\'{e}\ juxtapos\'{e}s,\ \}\idt }\par 
\vspace{\neglskip}\mbox{\blc  \blk\idt  \cmt  \{\ le\ robot\ partant\ du\ coin\ inf\'{e}rieur\ gauche\ de\ \ la\ \ ligne\ \ orient\'{e}\ \}\idt }\par 
\vspace{\neglskip}\mbox{\blc  \blk\idt  \cmt  \{\ vers\ l{'}est,\ et\ arrivant\ au\ coin\ inf\'{e}rieur\ droit\ \'{e}galement\ orien-\ \}\idt }\par 
\vspace{\neglskip}\mbox{\blc  \blk\idt  \cmt  \{\ t\'{e}\ \`{a}\ l{'}est.\ \ \ \ \ \ \ \ \ \ \ \ \ \ \ \ \ \ \ \ \ \ \ \ \ \ \ \ \ \ \ \ \ \ \ \ \ \ \ \ \ \ \ \ \ \ \ \ \ \ \ \ \ \ \}\idt }\par 
\pagebreak[3] 
\vspace{\neglskip}\mbox{\blc  \blk\blk\idt  \mcl var\idt \ I\ $:$\ integer$;$}\par 
\vspace{\neglskip}\mbox{\blc  \blk\idt  \mcl begin\idt }\par 
\pagebreak[3] 
\vspace{\neglskip}\mbox{\blc  \blk\blk\idt  \mcl for\idt \ I\ $:$$=$\ 1\ \mcl to\idt \ 7\ \mcl do\idt }\par 
\vspace{\neglskip}\mbox{\blc  \blk\blk\blk\idt  \mcl begin\idt \ }\par 
\vspace{\neglskip}\mbox{\blc  \blk\blk\blk\idt  \cmt  \{\ Le\ robot\ est\ crayon\ baiss\'{e},\ \`{a}\ l{'}est,\ en\ bas\ \`{a}\ gauche\ du\ motif\ \}\idt }\par 
\vspace{\neglskip}\mbox{\blc  \blk\blk\blk\blk\idt  Cercle$;$\ Carre$;$\ Triangles$;$\ }\par 
\vspace{\neglskip}\mbox{\blc  \blk\blk\blk\idt  \mcl end\idt $;$}\par 
\vspace{\neglskip}\mbox{\blc  \blk\idt  \mcl end\idt $;$}\par 
\medskip

\pagebreak[3] 
\vspace{\neglskip}\mbox{\blc  \blk\idt  \mcl procedure\idt \ AlaLigne$;$}\par 
\vspace{\neglskip}\mbox{\blc  \blk\blk\idt  \cmt  \{\ Ramener\ \ le\ \ robot\ \ du\ coin\ inf\'{e}rieur\ droit\ d{'}une\ ligne\ au\ coin\ \}\idt }\par 
\vspace{\neglskip}\mbox{\blc  \blk\blk\idt  \cmt  \{\ inf\'{e}rieur\ gauche\ de\ la\ ligne\ suivante\ sans\ \ changer\ \ l{'}orienta-\ \}\idt }\par 
\vspace{\neglskip}\mbox{\blc  \blk\blk\idt  \cmt  \{\ tion\ \ du\ \ robot\ \`{a}\ l{'}est,\ sans\ dessiner\ et\ en\ baissant\ le\ crayon\ \}\idt }\par 
\vspace{\neglskip}\mbox{\blc  \blk\blk\idt  \cmt  \{\ au\ d\'{e}but\ de\ la\ ligne\ suivante.\ \ \ \ \ \ \ \ \ \ \ \ \ \ \ \ \ \ \ \ \ \ \ \ \ \ \ \ \ \ \ \ \ \ \}\idt }\par 
\vspace{\neglskip}\mbox{\blc  \blk\idt  \mcl begin\idt }\par 
\vspace{\neglskip}\mbox{\blc  \blk\blk\idt  lc$;$\ ag$;$\ pqtd$;$\ av$;$\ av$;$\ pqtg$;$\ bc$;$}\par 
\vspace{\neglskip}\mbox{\blc  \blk\idt  \mcl end\idt $;$}\par 
\medskip

\pagebreak[3] 
\vspace{\neglskip}\mbox{\blc  \blk\idt  \mcl procedure\idt \ Pavage$;$}\par 
\vspace{\neglskip}\mbox{\blc  \blk\blk\idt  \cmt  \{\ Paver\ \ l{'}\'{e}cran\ \ de\ \ 11\ lignes\ commen\c{c}ant\ par\ un\ carr\'{e},\ puis\ par\ \}\idt }\par 
\vspace{\neglskip}\mbox{\blc  \blk\blk\idt  \cmt  \{\ des\ triangles\ puis\ par\ un\ cercle\ constitu\'{e}es\ de\ motifs\ \ r\'{e}p\'{e}ti-\ \}\idt }\par 
\vspace{\neglskip}\mbox{\blc  \blk\blk\idt  \cmt  \{\ tifs\ form\'{e}s\ d{'}un\ carr\'{e}\ suivi\ d{'}un\ cercle\ et\ de\ triangles\ juxta-\ \}\idt }\par 
\vspace{\neglskip}\mbox{\blc  \blk\blk\idt  \cmt  \{\ pos\'{e}s\ en\ partant\ du\ coin\ sup\'{e}rieur\ droit\ et\ en\ arrivant\ au\ coin\ \}\idt }\par 
\vspace{\neglskip}\mbox{\blc  \blk\blk\idt  \cmt  \{\ inf\'{e}rieur\ droit\ de\ l{'}\'{e}cran,\ avec\ le\ robot\ tourn\'{e}\ vers\ l{'}est.\ \ \ \ \}\idt }\par 
\pagebreak[3] 
\vspace{\neglskip}\mbox{\blc  \blk\blk\idt  \mcl var\idt \ J\ $:$\ integer$;$}\par 
\vspace{\neglskip}\mbox{\blc  \blk\idt  \mcl begin\idt }\par 
\vspace{\neglskip}\mbox{\blc  \blk\blk\idt  \cmt  \{\ Dessiner\ les\ lignes\ 1\ \`{a}\ 9\ de\ haut\ en\ bas.\ \ \ \ \ \ \ \ \ \ \ \ \ \ \ \ \ \ \ \ \ \ \ \}\idt }\par 
\pagebreak[3] 
\vspace{\neglskip}\mbox{\blc  \blk\blk\idt  \mcl for\idt \ J\ $:$$=$\ 1\ \mcl to\idt \ 3\ \mcl do\idt }\par 
\vspace{\neglskip}\mbox{\blc  \blk\blk\blk\idt  \mcl begin\idt }\par 
\vspace{\neglskip}\mbox{\blc  \blk\blk\blk\blk\idt  AlaLigne$;$\ LigneCarre$;$}\par 
\vspace{\neglskip}\mbox{\blc  \blk\blk\blk\blk\idt  AlaLigne$;$\ LigneTriangles$;$}\par 
\vspace{\neglskip}\mbox{\blc  \blk\blk\blk\blk\idt  AlaLigne$;$\ LigneCercle$;$}\par 
\vspace{\neglskip}\mbox{\blc  \blk\blk\blk\idt  \mcl end\idt $;$}\par 
\vspace{\neglskip}\mbox{\blc  \blk\blk\idt  \cmt  \{\ Dessiner\ les\ lignes\ 10\ puis\ 11.\ \ \ \ \ \ \ \ \ \ \ \ \ \ \ \ \ \ \ \ \ \ \ \ \ \ \ \ \ \ \ \ \ \}\idt }\par 
\vspace{\neglskip}\mbox{\blc  \blk\blk\idt  AlaLigne$;$\ LigneCarre$;$}\par 
\vspace{\neglskip}\mbox{\blc  \blk\blk\idt  AlaLigne$;$\ LigneTriangles$;$}\par 
\vspace{\neglskip}\mbox{\blc  \idt  \mcl end\idt $;$}\par 
\medskip

\vspace{\neglskip}\mbox{\blc  \idt  \mcl begin\idt }\par 
\vspace{\neglskip}\mbox{\blc  \blk\idt  \cmt  \{\ Placer\ le\ robot\ au\ coin\ sup\'{e}rieur\ droit\ orient\'{e}\ vers\ l{'}est.\ \ \ \ \ \ \}\idt }\par 
\vspace{\neglskip}\mbox{\blc  \blk\idt  eh$;$\ ad$;$\ pqtd$;$}\par 
\vspace{\neglskip}\mbox{\blc  \blk\idt  \cmt  \{\ Paver\ l{'}\'{e}cran,\ terminer\ au\ coin\ inf\'{e}rieur\ droit\ orient\'{e}\ \`{a}\ l{'}est.\ \}\idt }\par 
\vspace{\neglskip}\mbox{\blc  \blk\idt  Pavage$;$\ st$;$\ }\par 
\vspace{\neglskip}\mbox{\blc  \idt  \mcl end\idt $.$}\par 

\medskip }


\noindent $\Box$}\end{corriges}

%%%%%%%%%%%%%%%%%%%%%%%%%%%%%%%%%%%%%%%%%%%%%%%%%%%%%%%%%%%%%%%%%%%%%%%%%%%
\noindent \framebox[\textwidth]{
\begin{minipage}{\myparboxwidth}
\setlength{\parindent}{2em}
\small \sf
\vspace{\abovedisplayskip}
\section*{Addenda}

La boucle ``for'' est une des premières notions linguistiques qui soit
apparue en informatique puisqu'elle fût inventée à l'ETH de Zürich
 en 1952 par Heinz Rutishauser \cite{rutishauser}.

La comparaison des dessins que l'on peut obtenir en utilisant la boucle
``for'' avec ceux que l'on peut faire en mode interactif donne la mesure
de la puissance des ordinateurs quand il s'agit d'effectuer rapidement des
tâches répétitives. En cela l'ordinateur se distingue très
nettement des calculettes non programmables. On comprend également
que cette puissance de calcul n'a rien à voir avec une quelconque forme
de pensée ou d'intelligence si ce n'est celle du programmeur de
l'ordinateur.

En \pascal, on utilise une boucle {\mcl for} {\idt I} $:=$ $n_1$ {\mcl to}
 $n_2$ {\mcl do} {\mcl begin} \ldots {\mcl end}$;$ pour répéter
$(n_2 - n_1) + 1$ fois le corps de la boucle compris entre {\mcl begin} et
{\mcl end} sauf si $n_2 < n_1$ auquel cas le corps de la boucle n'est pas 
exécuté. A chaque itération, le compteur de boucle {\idt I} prend les
valeurs successives $n_1$, $n_1 + 1$, $n_1 + 2$, \ldots, $n_2$. Dans la forme {\mcl for} 
{\idt I} $:=$ $n_1$ {\mcl downto} $n_2$ {\mcl do} {\mcl begin} \ldots {\mcl end}$;$,
{\idt I} prend les valeurs successives $n_1$, $n_1 - 1$, $n_1 - 2$, \ldots, $n_2$
sauf si $n_2 > n_1$ auquel cas la boucle n'est pas exécutée (le mot anglais
{\mcl downto} signifie `en descendant jusqu'à').

L'utilisation de boucles ``for'' exerce la capacité de percevoir
quand un dessin peut se décomposer en un certain nombre de sous-dessins
identiques (dans le cas d'une boucle simple) ou similaires (dans le cas
de boucles imbriquées). C'est un exercice quotidien pour l'informaticien
qui doit savoir décomposer une tâche complexe en sous-tâches simples,
ce qui demande parfois énormément d'imagination.
Il s'agit d'un premier pas vers l'acquisition du raisonnement
inductif des scien\-ti\-fi\-ques qui, à partir d'observations spécifiques,
créent une hypothèse générale rendant compte de ces observations,
hypothèse qu'il reste à valider dans tous les cas possibles, ici de 
manière tout à fait expérimentale, en observant le dessin à 
l'écran. Les mathématiciens appellent {\em conjecture} une telle
hypothèse induite de l'expérience et non encore démontrée.

Ce raisonnement inductif a longtemps été confondu avec le raisonnement 
par récurrence qui est essentiel en mathématiques. Il a fallu plus d'un 
millénaire pour que les mathématiciens
(Fermat\footnote{\ \href{https://en.wikipedia.org/wiki/Pierre_de_Fermat}{Pierre de Fermat}, 1601--1665, mathématicien français.}, 
Pascal\footnote{\ \href{https://en.wikipedia.org/wiki/Blaise_Pascal}{Blaise Pascal}, 1623--1662, mathématicien français.}
notamment) comprennent
pleinement ce raisonnement par récurrence (\cite{knuth2} page 17).
\vspace{\abovedisplayskip}
\end{minipage}
}
%%%%%%%%%%%%%%%%%%%%%%%%%%%%%%%%%%%%%%%%%%%%%%%%%%%%%%%%%%%%%%%%%%%%%%%%%%%%

%%%%%%%%%%%%%%%%%%%%%%%%%%%%%%%%%%%%%%%%%%%%%%%%%%%%%%%%%%%%%%%%%%%%%%%%%%%
\noindent \framebox[\textwidth]{
\begin{minipage}{\myparboxwidth}
\setlength{\parindent}{2em}
\small \sf
\vspace{\abovedisplayskip}

 Nous en
sommes pour l'instant aux prémisses, l'absence d'induction correcte 
 étant encore admise dans le raisonnement
par récurrence~: si le dessin est bien 
fait pour le motif de base et pour
la première répétition du motif alors il est réputé correct, seul le 
problème du cadrage dans la fenêtre restant à résoudre par simple
expérimentation.

La notion d'{\em invariant} (c'est-à-dire pour notre \robot\ la
description de ses positions et orientations possibles au début du
corps des boucles) est due à Turing
\footnote{\ \href{https://en.wikipedia.org/wiki/Alan_Turing}{Alan Turing}, informaticien anglais.} en 1949, oubliée, puis redécouverte par
Naur%
\footnote{\ \href{https://en.wikipedia.org/wiki/Peter_Naur}{Peter Naur}, informaticien danois.} en 1966 et Floyd%
\footnote{\ \href{https://en.wikipedia.org/wiki/Robert_W._Floyd}{Robert Floyd}, informaticien américain.} en 1967. 
Elle repose
sur un raisonnement par récurrence~: si l'invariant est vrai quand on
arrive pour la première fois dans la boucle (c'est-à-dire que la
position et l'orientation du \robot\ est conforme à la description qui en
est faite dans l'invariant) et si, en supposant que l'invariant soit vrai au
début de l'exécution du corps de la boucle, il reste encore vrai après
l'exécution de ce corps de boucle alors l'invariant est toujours vrai au 
début de chaque itération dans la boucle et l'est encore quand
on sort de la boucle. La notion s'applique aussi bien aux procédures~: on 
peut utiliser un 
{\em invariant d'entrée} décrivant les positions et orientations du
\robot\ avant l'appel et un {\em invariant de sortie} décrivant 
ces positions et orientations après l'appel, de sorte que le corps de la
procédure soit compréhensible sans avoir à consulter tous les points
d'appel pour savoir dans quelles conditions la procédure est appelée.

L'ordinateur peut être utilisé pour manipuler des programmes, en
particulier pour découvrir automatiquement des invariants de forme
simple en chaque point d'un programme. C'est l'objet de 
l'{\em interprétation abstraite} des programmes
\cite{CousotIA}. Pour le \robot, par l'exemple, une forme simple
d'invariant est constituée par l'ensemble de ses orientations possibles, 
en ignorant la position de son crayon, ses coordonnées sur 
l'écran, etc. Les orientations possibles du \robot\ concret peuvent être
calculées en exécutant le programme avec un \robot\ abstrait qui
exécute les commandes du \robot\ concret sans bouger
mais en changeant simplement d'orientation. Pour ce \robot\ abstrait
la commande {\idt av} ne fait rien car elle ne change pas l'orientation du
\robot\ concret$\:$; la commande {\idt vd} a le même effet que 
{\idt pqtd}$\:$; etc. Pendant l'exécution abstraite les orientations 
possibles du \robot\ abstrait sont notées comme commentaire
en chaque point du programme. L'exécution abstraite du programme
s'arrête après stabilisation des résultats. Par exemple, pour le programme {\idt FriseDeN\_2} (page 
\pageref{prog:FriseDeN_2}), on trouve~:
\vspace{\abovedisplayskip}
\end{minipage}
}
%%%%%%%%%%%%%%%%%%%%%%%%%%%%%%%%%%%%%%%%%%%%%%%%%%%%%%%%%%%%%%%%%%%%%%%%%%%%

%%%%%%%%%%%%%%%%%%%%%%%%%%%%%%%%%%%%%%%%%%%%%%%%%%%%%%%%%%%%%%%%%%%%%%%%%%%
\noindent \framebox[\textwidth]{
\begin{minipage}{\myparboxwidth}
\setlength{\parindent}{2em}
\small \sf
\vspace{\abovedisplayskip}


{
\parindent \parindentprog
\medskip

\carrenoir
\vspace{\neglskip}\mbox{\blc  \idt  \mcl program\idt \ FriseDeN\_2$;$
}\par 
\vspace{\neglskip}\mbox{\blc \blk\idt \mcl uses\idt \ \ldots\ $;$
}\par 
\pagebreak[3] 
\vspace{\neglskip}\mbox{\blc  \blk\idt  \mcl var\idt \ J\ $:$\ integer$;$}\par 
\vspace{\neglskip}\mbox{\blc  \idt  \mcl begin\idt \ }\par 
\vspace{\neglskip}\mbox{\blc  \blk\idt  \cmt  \{\ Nord\ \}\idt }\par 
\vspace{\neglskip}\mbox{\blc  \blk\idt  ag$;$}\par 
\vspace{\neglskip}\mbox{\blc  \blk\idt  \cmt  \{\ Nord\ \}\idt }\par 
\pagebreak[3] 
\vspace{\neglskip}\mbox{\blc  \blk\idt  \mcl for\idt \ J\ $:$$=$\ 1\ \mcl to\idt \ 16\ \mcl do\idt }\par 
\vspace{\neglskip}\mbox{\blc  \blk\blk\idt  \mcl begin\idt }\par 
\vspace{\neglskip}\mbox{\blc  \blk\blk\blk\idt  \cmt  \{\ Nord\ \}\idt \ av$;$\ \cmt  \{\ Nord\ \}\idt \ av$;$\ \cmt  \{\ Nord\ \}\idt \ pdt$;$\ \cmt  \{\ Sud\ \}\idt \ }\par 
\vspace{\neglskip}\mbox{\blc  \blk\blk\blk\idt  vg$;$\ \cmt  \{\ Est\ \}\idt \ vd$;$\ \cmt  \{\ Sud\ \}\idt \ pdt$;$\ \cmt  \{\ Nord\ \}\idt \ }\par 
\vspace{\neglskip}\mbox{\blc  \blk\blk\idt  \mcl end\idt $;$}\par 
\vspace{\neglskip}\mbox{\blc  \blk\idt  \cmt  \{\ Nord\ \}\idt \ }\par 
\vspace{\neglskip}\mbox{\blc  \blk\idt  st$;$}\par 
\vspace{\neglskip}\mbox{\blc  \blk\idt  \cmt  \{\ Nord\ \}\idt \ }\par 
\vspace{\neglskip}\mbox{\blc  \idt  \mcl end\idt $.$}\par 

\medskip }



La règle des signes est un exemple d'interprétation abstraite utilisée
en mathématiques. Le signe d'une expression peut être
calculé en utilisant le signe des variables, qui remplace leurs valeurs,
grâce aux opérations abstraites comme $\oplus + \oplus = \oplus$, 
$\oplus - \ominus = \oplus$, $\oplus \times \ominus = \ominus$, etc. Le prix
à payer pour cette simplification est évidemment l'imprécision des
résultats, comme pour $\oplus + \ominus$.
\vspace{\abovedisplayskip}
\end{minipage}
}
%%%%%%%%%%%%%%%%%%%%%%%%%%%%%%%%%%%%%%%%%%%%%%%%%%%%%%%%%%%%%%%%%%%%%%%%%%%%

